\section*{\textit {Abstract}} %%%% (DO NOT CHANGE) %%%%

This study addresses a fundamental question in developing statistical monitoring systems: can SAR observations from the same class but different areas be modeled by a common distribution? 
Motivated by increasing anthropogenic pressures on the Cerrado and the growing use of SAR data for change detection, we investigate the spatial behavior of sample distributions for the Cerrado’s forest, savanna, and grassland formations using NISAR and Sentinel-1 imagery. 
We compare the goodness of fit of the $\Gamma_{\text{SAR}}$ and $\mathcal{G}^0_I$ models and assess pairwise similarity using Kullback–Leibler-based hypothesis tests. 
The $\mathcal{G}^0_I$ model provides the best fit across all classes, polarizations, and missions. 
However, pairwise tests reveal substantial within-class differences associated with mean brightness, indicating that a single set of parameters is inadequate to describe the spatial distribution of observations. 
Likelihood-ratio tests support sharing shape parameter $\alpha$ while allowing $\mu$ to vary, except for NISAR Grassland (HH/HV).
These findings suggest representing classes as families of $\mathcal{G}^0_I$ distributions.
\\

%%%% Add the Key words below %%%%
\textit{\textbf{Key words --} %%%% (DO NOT CHANGE) %%%%
SAR, Cerrado, Sentinel-1, NISAR, Change Detection, Statistical Analysis, Stochastic Distances. %%%% (ADD KEY WORDS) %%%%
}